\section*{Chapitre \ref{ch:m1:stock-loan-and-short-selling} --- Prêt de titres et vente à découvert en pratique}

\begin{solution}{exo:m1:stock-loan-and-short-selling:1}
\omterm{def:m1:financing:seclending}{Frais d'emprunt} $= 4.3\% - (-11.7\%) = 16\%$. Coût quotidien $3\,000\,000 \times 0.16/360 =
\$1\,333$.
\end{solution}

\begin{solution}{exo:m1:stock-loan-and-short-selling:2}
\omterm{def:m1:stock-loan-and-short-selling:utilisation}{Taux d'utilisation} $10.2/12 = 85\%$. Sur la courbe~: $0.3\% + 79.7\% \times
\bigl((0.85-0.75)/0.25\bigr)^2 = 0.3\% + 79.7\% \times 0.16 = 13.05\%$.
\end{solution}

\begin{solution}{exo:m1:stock-loan-and-short-selling:3}
30~\% du flottant~; 7.5~\omterm{def:m1:stock-loan-and-short-selling:si}{jours de couverture}. Les \omterm{def:m1:stock-loan-and-short-selling:si}{jours de couverture} comptent davantage pour
le \omterm{def:m1:stock-loan-and-short-selling:squeeze}{short squeeze} lui-même~: ils mesurent le temps que mettraient les vendeurs à découvert, tous acheteurs, pour
sortir, d'où l'impact $k$ de leur couverture. L'\omterm{def:m1:stock-loan-and-short-selling:si}{encours vendeur à découvert} rapporté au flottant
indique à quel point l'emprunt est tendu, et prédit donc les \omterm{def:m1:financing:seclending}{frais d'emprunt} et les \omterm{def:m1:stock-loan-and-short-selling:recall}{rappels}.
\end{solution}

\begin{solution}{exo:m1:stock-loan-and-short-selling:4}
$0.25/0.40 = 0.625$ an $= 225$ jours~; $0.25/0.80 = 112$ jours. La date
de la chute, répartie uniformément sur 720 jours, seuls les scénarios où elle tombe
avant le jour 225 sont gagnants~: environ 31~\%. Une position qui a raison avec certitude sur la
destination perd dans deux scénarios sur trois.
\end{solution}

\begin{solution}{exo:m1:stock-loan-and-short-selling:5}
A prête 100 actions à S$_1$, qui les vend à B. Le courtier de B prête les 100 à
S$_2$, qui les vend à C. L'\omterm{def:m1:stock-loan-and-short-selling:si}{encours vendeur à découvert} est de 200. Pour 123~: S$_2$ emprunte et
vend seulement 23 actions de B. Toutes les ventes ont été livrées. Détenteurs longs~: A dispose
d'une créance sur 100, B détient 100 actions dont 23 sont prêtées, C en détient 23~: positions
longues économiques de 223 actions, soit le flottant plus l'\omterm{def:m1:stock-loan-and-short-selling:si}{encours vendeur à découvert}. Seuls les propriétaires légaux
à la \omterm{def:m1:shares-corporate-actions-indices:dividend}{date d'enregistrement} votent~: B pour 77 actions et C pour 23~; A et B ne peuvent pas voter les
actions qu'ils ont prêtées, à moins de les rappeler.
\end{solution}

\begin{solution}{exo:m1:stock-loan-and-short-selling:6}
Date de règlement-livraison~: mardi 15 septembre. Close-out avant le début de la séance
régulière du jour de règlement suivant~: l'ouverture du mercredi 16
septembre. En cas de vente longue en défaut de livraison~: le troisième jour de règlement suivant la
date de règlement, l'ouverture du vendredi 18 septembre.
\end{solution}

\begin{solution}{exo:m1:stock-loan-and-short-selling:7}
P\&L net maximal~: 789\,000~\$ au jour 190. P\&L de prix final~: $+\$164\,000$. Coût d'emprunt
cumulé final~: 399\,000~\$, dont 308\,000~\$, soit 77~\%, après le jour
190. Net final~: $-\$235\,000$.
\end{solution}

\begin{solution}{exo:m1:stock-loan-and-short-selling:8}
L'\omterm{def:m1:stock-loan-and-short-selling:si}{encours vendeur à découvert} compte les positions courtes, et non les actions manquantes. Chaque vente à découvert
a été réglée avec une action empruntée~; son acheteur est devenu pleinement propriétaire, et son
courtier peut prêter de nouveau cette action. La chaîne peut être parcourue plus d'une fois,
donc un encours vendeur supérieur à 100~\% relève de l'arithmétique, comme le montre le rapport des services du
régulateur sur le début de 2021 l'explique avec le même exemple. La preuve d'une \omterm{def:m1:stock-loan-and-short-selling:naked}{vente à découvert à nu} est
un \emph{défaut de livraison}~: des défauts persistants à la chambre de compensation, la
présence du titre sur la liste des seuils, les close-out au titre de la Rule 204. Ces données sont
publiées, et forment un jeu de données distinct de l'\omterm{def:m1:stock-loan-and-short-selling:si}{encours vendeur à découvert}.
\end{solution}

\begin{solution}{pb:m1:stock-loan-and-short-selling:1}
\textbf{1.} 40~\% du flottant~; 5~\omterm{def:m1:stock-loan-and-short-selling:si}{jours de couverture}~; le fonds représente 10~\% de l'\omterm{def:m1:stock-loan-and-short-selling:si}{encours
vendeur à découvert}.
\textbf{2.} $40\,000\,000 \times 0.12/360 = \$13\,333$.
\textbf{3.} $0.30/0.12 = 2.5$ ans~: 900 jours.
\textbf{4.} 66\,667~\$ par jour~; 180 jours.
\textbf{5.} La position moyenne dure 270 jours et les \omterm{def:m1:financing:seclending}{frais d'emprunt} absorbent désormais la
totalité du gain espéré en 180 jours~: à ce niveau de frais, la position a une valeur
espérée négative, à moins que le fonds n'ait une raison d'anticiper la baisse plus tôt que d'habitude. La
décision de rester vendeur à découvert est une nouvelle décision, à prendre au nouveau prix de
l'emprunt.
\textbf{6.} $x = (0.25 - 0.2 \times 0.1/0.5)/(1 - 0.2/0.5) = 0.21/0.6 = 35\%$;
la moitié des vendeurs à découvert se sont couverts.
\textbf{7.} La formule donne $(0.25 - 0.08)/0.2 = 85\% > b$~: chaque vendeur à découvert
se couvre et $x = 25\% + 40\% = 65\%$.
\textbf{8.} $x$ atteint 60~\% lorsque $(s - 0.08)/0.2 = 0.6$~: $s = 20\%$.
\textbf{9.} Cours de 33~\$~: $2\,000\,000 \times 13 = 26$ millions de dollars, soit 65~\% de la
valeur initiale de la position.
\textbf{10.} Valeur de marché de 66 millions de dollars~; la marge passe de 30~\% à 100~\%~:
46.2 millions de dollars de plus, le jour où la position a perdu 26 millions de dollars.
\textbf{11.} $500\,000 \times 13 = 6.5$ millions de dollars.
\textbf{12.} Les prêteurs sont des détenteurs longs~: une hausse de 65~\% est le moment où ils vendent, et une
vente est un \omterm{def:m1:stock-loan-and-short-selling:recall}{rappel}. Certains rappellent pour prêter à nouveau, à des frais plus élevés. Et une
OPA suppose un vote, pour lequel les actions doivent être rappelées.
\textbf{13.} $40\,000\,000 \times 20\% \times 90/360 = 2$ millions de dollars, contre
1.2 millions de dollars au taux variable de 12~\%. Les 0.8 million de dollars supplémentaires achètent trois mois
sans \omterm{def:m1:stock-loan-and-short-selling:recall}{rappel} ni révision du prix~: une assurance contre précisément les parties II et
III, peu chère à côté d'un \omterm{def:m1:stock-loan-and-short-selling:recall}{buy-in} de 6.5 millions de dollars.
\textbf{14.} Acheter des puts, ou un put spread pour réduire la prime. La perte est
limitée par la prime, il n'y a ni \omterm{def:m1:stock-loan-and-short-selling:recall}{rappel} ni spirale de marge. Les \omterm{def:m1:financing:seclending}{frais
d'emprunt} n'ont pas disparu~: le \omterm{def:m1:what-a-trading-firm-does:market-maker}{teneur de marché} qui vend les puts se couvre en
vendant le titre à découvert et en supporte le coût~; celui-ci est donc inclus dans le prix du put, comme un prix à terme
inférieur au cours au comptant (question d'entretien 6).
\textbf{15.} \omterm{def:m1:stock-loan-and-short-selling:si}{Encours vendeur à découvert} de 40~\% du flottant, avec cinq \omterm{def:m1:stock-loan-and-short-selling:si}{jours de couverture}~:
encombré et lent à sortir. La part de 10~\% de l'\omterm{def:m1:stock-loan-and-short-selling:si}{encours vendeur à découvert} que détient le fonds
signifie qu'il ne peut pas sortir sans faire bouger le cours.
\textbf{16.} Le rapport des services du régulateur sur janvier 2021 constate que les couvertures par les vendeurs à
découvert ne représentaient qu'une petite fraction des achats, et que c'est le sentiment, et non les couvertures,
qui a soutenu la hausse. Une forte hausse d'un titre très vendu à découvert est ce que produit la
cascade décrite à la \cref{prop:m1:stock-loan-and-short-selling:cascade} à partir
d'un choc ordinaire~; la manipulation est une affirmation sur l'intention et la coordination.
\textbf{17.} Juste sur la valeur~: passer de 20~\$ à 13~\$ représente une baisse de 35~\%, soit 14
millions de dollars sur 2 millions d'actions. Il n'en a rien encaissé sur les 500\,000 actions
achetées à 33~\$, et n'aurait encaissé le reste que s'il avait pu faire face à un
\omterm{def:m1:financing:margincall}{appel de marge} de 46 millions de dollars. Avoir raison sur la destination est une
condition nécessaire, mais loin d'être suffisante, pour une position vendeuse.
\textbf{18.} Pour le scénario de \omterm{def:m1:stock-loan-and-short-selling:squeeze}{short squeeze}, et non celui de la volatilité~: dimensionner la position de telle sorte qu'un
mouvement jusqu'au sommet de la fourchette de couverture ($+65\%$ ici) et la hausse de marge qui y est associée
puissent être financés sans vendre autre chose~; limiter la part
de l'\omterm{def:m1:stock-loan-and-short-selling:si}{encours vendeur à découvert} de tout titre et les jours de volume détenus~; privilégier l'emprunt à terme
ou des options lorsque l'\omterm{def:m1:stock-loan-and-short-selling:si}{encours vendeur à découvert} est élevé.
\textbf{19.} \textbf{180 jours.}
\textbf{20.} L'action elle-même, auprès de détenteurs de long terme qui peuvent la réclamer
à tout moment, à un coût qui augmente lorsque la position est prisée, et lorsqu'elle perd
de l'argent.
\end{solution}

\begin{solution}{iq:m1:stock-loan-and-short-selling:1}
Le fonds demande à son \omterm{def:m1:financing:pb}{prime broker} un locate~; le courtier confirme qu'il peut
se procurer les actions. Le fonds vend sur le marché. Pour le règlement-livraison, le courtier
emprunte les actions, auprès des clients qui les ont déposées en marge ou auprès d'un prêteur
agent, en déposant un collatéral en espèces légèrement supérieur à leur valeur, et
les livre à l'acheteur. Chaque jour, le prêt fait l'objet d'une \omterm{def:m1:pnl-and-positions:mtm}{valorisation au prix de marché}, le fonds
verse les \omterm{def:m1:financing:seclending}{frais d'emprunt} (ou perçoit le \omterm{def:m1:financing:seclending}{taux de rebate}) et doit tout dividende éventuel. Pour clôturer, le fonds
rachète les actions~; le courtier les restitue et récupère son collatéral. À tout
moment, le prêteur peut rappeler les titres, et les frais peuvent changer.
\iqlookfor{locate avant la vente, le collatéral, la valorisation quotidienne, le risque de \omterm{def:m1:stock-loan-and-short-selling:recall}{rappel}.}
\end{solution}

\begin{solution}{iq:m1:stock-loan-and-short-selling:2}
L'acheteur d'une action vendue à découvert en est pleinement propriétaire et peut la prêter de nouveau~; la même
action peut servir de support à plusieurs positions vendeuses. Le total des positions longues est égal au
flottant augmenté de l'\omterm{def:m1:stock-loan-and-short-selling:si}{encours vendeur à découvert}. Cette égalité ne suppose aucun échec de livraison.
\iqlookfor{la chaîne de re-prêt en deux phrases, et l'identité longs
$=$ flottant $+$ positions vendeuses.}
\end{solution}

\begin{solution}{iq:m1:stock-loan-and-short-selling:3}
Coûts d'emprunt et disponibilité. Le décile inférieur de nombreux signaux est
peuplé de petites valeurs, fortement vendues à découvert~: des frais de dizaines de pourcent peuvent
dépasser les 9~\%, certaines valeurs ne peuvent pas être localisées du tout, et les positions sont
rappelées aux pires moments. Le backtest exige des données historiques de frais et
de disponibilité, et une règle pour les \omterm{def:m1:stock-loan-and-short-selling:recall}{buy-ins}. Vérifier aussi si le rendement de la jambe
vendeuse résiste à l'exclusion des \omterm{def:m1:stock-loan-and-short-selling:special}{special}~: souvent, il s'agit des frais déguisés.
\iqlookfor{l'idée que l'anomalie peut \emph{être} les \omterm{def:m1:financing:seclending}{frais d'emprunt}.}
\end{solution}

\begin{solution}{iq:m1:stock-loan-and-short-selling:4}
Cette demande d'emprunt a dépassé l'\omterm{def:m1:stock-loan-and-short-selling:utilisation}{offre prêtable}~: soit de nombreuses nouvelles ventes à découvert,
souvent informées, ou un
retrait de l'offre, par exemple des détenteurs qui rappellent leurs titres avant un vote ou
vendent. Elle ne dit pas lequel, ni que le titre baissera assez vite pour
payer 45~\%. Elle dit en revanche qu'un \omterm{def:m1:stock-loan-and-short-selling:squeeze}{short squeeze} et des \omterm{def:m1:stock-loan-and-short-selling:recall}{rappels} sont devenus plus probables,
et que les prix des options vont se déplacer pour refléter les frais.
\iqlookfor{demande contre offre, et la distinction entre information et
calendrier.}
\end{solution}

\begin{solution}{iq:m1:stock-loan-and-short-selling:5}
Un \omterm{def:m1:what-a-trading-firm-does:market-maker}{teneur de marché} doit vendre sur demande, y compris des actions qu'il ne détient pas, pour
maintenir une cotation bilatérale~; un locate avant chaque vente est impossible à cette
vitesse. L'exception est encadrée~: elle ne couvre que le market making de bonne foi, et non
les positions spéculatives ni la cotation unilatérale~; les fails doivent néanmoins être clôturés,
dans un délai légèrement plus long~; et les régulateurs examinent les sociétés dont l'activité ne
ressemble pas au market making.
\iqlookfor{pourquoi l'exception existe et que l'obligation de clôture
subsiste.}
\end{solution}

\begin{solution}{iq:m1:stock-loan-and-short-selling:6}
Non. La parité call-put est $C - P = S - K\mathrm{e}^{-rT}$ seulement si le titre peut être vendu à découvert
sans aucun coût. Pour exploiter des put «~chers~», on vend le put, on achète le call
et vend à découvert le titre~: la position vendeuse paie les \omterm{def:m1:financing:seclending}{frais d'emprunt} $\phi$. La bonne
forward est $S\mathrm{e}^{(r-\phi)T}$, et la violation apparente traduit
les \omterm{def:m1:financing:seclending}{frais d'emprunt} implicites du marché, souvent un meilleur indicateur des frais que
toute cotation d'un courtier. Les risques restants sont le \omterm{def:m1:stock-loan-and-short-selling:recall}{rappel} et une variation des frais, ce qui
explique que les frais implicites puissent rester au-dessus de ceux affichés.
\iqlookfor{les frais comme rendement de type dividende dans le forward~; emprunt implicite.}
\end{solution}
